\section{State-of-the-Art in Anomaly Detection}
\label{sec:sota_anomaly_detection}

Time-series anomaly detection concerns the identification of observations, subsequences, or events that deviate from expected behaviour. Unlike anomaly detection in i.i.d.\ data, temporal abnormality may depend on seasonality, exogenous context, multivariate relationships, and changing uncertainty \cite{chandola2009anomaly,blazquez2021review}. These characteristics are particularly relevant in electricity markets and grid operations, where weather, calendar conditions, and market state influence the interpretation of unusual patterns \cite{hong2016probabilistic,weron2014electricity}. This section reviews the principal deep-learning paradigms used for time-series anomaly detection and their suitability for operational energy applications.

\subsection{What is Anomaly Detection? Theoretical Framing}

Let $\mathbf{x}_t \in \mathbb{R}^{d}$ denote a multivariate observation at time $t$, and let $\mathbf{X}_{1:T} = \{\mathbf{x}_1,\mathbf{x}_2,\dots,\mathbf{x}_T\}$ denote the observed series. A general anomaly detector constructs an anomaly score
\begin{equation}
	s_t = f_{\theta}\!\left(\mathbf{x}_{1:t}, \mathbf{c}_t\right),
	\label{eq:sota_score}
\end{equation}
where $\mathbf{c}_t$ denotes contextual information such as time-of-day, holiday state, weather conditions, regime information, or latent historical state, and $f_{\theta}(\cdot)$ is a learned or designed scoring function. A binary anomaly decision is then produced through thresholding,
\begin{equation}
	z_t = \mathbb{I}(s_t > \tau_t),
	\label{eq:sota_binary}
\end{equation}
where $\tau_t$ may be fixed, learned from validation data, or conditioned on context. This abstract formulation covers most anomaly detection paradigms in the literature. Reconstruction-based methods derive $s_t$ from reconstruction error; forecasting-based methods derive it from prediction residuals; density-based methods derive it from likelihood or energy; and attention-based methods derive it from abnormal association patterns in latent representations \cite{chandola2009anomaly,blazquez2021review}.

In operational time-series domains, anomaly detection is rarely a fully supervised classification task because clean anomaly labels are scarce, incomplete, or disputed. Instead, most systems are unsupervised or semi-supervised: they learn a compact model of normality and then treat significant departures from that learned normality as anomalies. This is especially true in electricity markets and power systems, where abnormalities may reflect rare market regimes, curtailment, outages, telemetry faults, weather-induced variability, or cross-variable inconsistencies, but where a single agreed label is often unavailable.

\subsection{Types of Anomalies in Time Series}

The theoretical background for the later methodology requires distinguishing several anomaly types.

\textbf{Point anomalies} are isolated observations whose values are individually abnormal relative to recent history or expected behaviour. In energy systems, a one-hour price spike or a sudden load drop may fall into this category.

\textbf{Contextual anomalies} are observations that are anomalous only under a specific context. A low photovoltaic output at night is normal, whereas the same output at midday under high solar radiation is suspicious. Such anomalies are central in weather-driven power systems.

\textbf{Collective anomalies} arise when a sequence of points is anomalous as a group even if each point alone is not sufficiently extreme. Prolonged volatility, sustained bias, or regime shifts are typical examples.

\textbf{Multivariate relationship anomalies} occur when the joint relationship among variables becomes abnormal even though each variable remains plausible in isolation. For example, electricity price may remain unusually low despite high load and weak renewable generation. In this case, the anomaly lies in the cross-variable relationship rather than in a marginal outlier.

These categories are methodologically important because they determine what type of detector is appropriate. Point anomalies may be detected by local scoring; contextual anomalies require conditioning on exogenous variables; collective anomalies require sequence-level reasoning; and multivariate relationship anomalies require models capable of learning joint temporal structure rather than independent scalar thresholds \cite{chandola2009anomaly,blazquez2021review}.

\subsection{Core Paradigms in Deep Time-Series Anomaly Detection}

Modern deep anomaly detectors usually follow one of four broad paradigms.

\textbf{Reconstruction-based detection} trains a model to compress and reconstruct normal data. Large reconstruction error suggests abnormality:
\begin{equation}
	s_t^{(\mathrm{rec})} = \|\mathbf{x}_t - \hat{\mathbf{x}}_t\|_2^2.
	\label{eq:sota_reconstruction}
\end{equation}
Autoencoders, convolutional autoencoders, recurrent autoencoders, variational autoencoders, and memory-augmented autoencoders belong to this family.

\textbf{Forecasting-based detection} learns to predict a future target or sequence and flags observations when prediction residuals are unusually large:
\begin{equation}
	s_t^{(\mathrm{pred})} = \|\mathbf{y}_t - \hat{\mathbf{y}}_t\|_2.
	\label{eq:sota_forecast}
\end{equation}
This paradigm is especially suitable for weakly labelled operational settings because it directly models normal dynamics.

\textbf{Density- or probability-based detection} estimates the probability of observations or latent states and identifies low-likelihood regions as anomalous. Variational recurrent models and deep Gaussian mixture approaches are representative examples.

\textbf{Association-discrepancy detection} is typical of recent Transformer-based methods. Instead of relying only on residual magnitude, such methods define anomalies through abnormal internal dependency patterns or mismatched attention associations \cite{xu2022anomalytransformer,tuli2022tranad}.

From the viewpoint of this thesis, forecasting-based detection is particularly relevant because it aligns naturally with multivariate energy monitoring under weak labels. Nevertheless, autoencoders, CNNs, LSTMs, and Transformers all contribute important ideas to the state of the art and therefore deserve critical review.

\subsection{Autoencoder-Based Detectors}

Autoencoders are among the most influential foundations of deep anomaly detection. They learn an encoder-decoder mapping
\begin{equation}
	\mathbf{h}_t = f_{\psi}(\mathbf{x}_t), \qquad \hat{\mathbf{x}}_t = g_{\phi}(\mathbf{h}_t),
	\label{eq:sota_ae}
\end{equation}
where the encoder $f_{\psi}(\cdot)$ compresses the input into a latent representation and the decoder $g_{\phi}(\cdot)$ reconstructs it. The anomaly score is usually derived from reconstruction error, as in \eqref{eq:sota_reconstruction}. The underlying assumption is that the autoencoder learns a compact manifold of normal behaviour and therefore reconstructs normal samples well while reconstructing abnormal samples poorly.

Autoencoder-based methods are attractive because they do not require labels and naturally support multivariate modeling. Early work demonstrated that nonlinear autoencoders can detect anomalies more flexibly than linear dimensionality-reduction methods \cite{zhou2017robustae}. Subsequent approaches expanded this idea in several directions. Deep Autoencoding Gaussian Mixture Models (DAGMM) combined autoencoding with probabilistic density estimation in latent space to improve detection of low-density events \cite{zong2018dagmm}. Memory-augmented autoencoders explicitly enforced a memory of normal patterns to sharpen separation between normal and abnormal samples \cite{gong2019memorizing}. USAD proposed an unsupervised adversarial autoencoder-style architecture tailored to multivariate time-series anomaly detection \cite{audibert2020usad}. Variational and sequence-aware autoencoders further broadened the design space by combining latent uncertainty modeling with temporal structure \cite{su2019omnianomaly}.

The strengths of autoencoders are clear. They are flexible, unsupervised, and naturally suitable when normal data dominate the training set. They also offer a direct and interpretable anomaly mechanism through reconstruction error. However, their limitations are equally important. First, highly expressive autoencoders can sometimes reconstruct anomalies too well, thereby reducing the contrast between normal and abnormal patterns. Second, reconstruction error alone does not necessarily tell whether an anomaly is operationally meaningful, physically plausible, or economically relevant. Third, generic autoencoders often ignore explicit domain rules unless those rules are incorporated externally.

\subsection{CNN-Based Time-Series Anomaly Detectors}

Convolutional neural networks became important in time-series anomaly detection because one-dimensional convolutions can learn local motifs, shift-invariant temporal patterns, and multi-scale structures efficiently. A temporal convolution layer applies a kernel over recent observations:
\begin{equation}
	h_t^{(k)} = \sigma\!\left(\sum_{i=0}^{K-1} w_i^{(k)} x_{t-i} + b^{(k)} \right),
	\label{eq:sota_cnn}
\end{equation}
where $K$ is the kernel size and $w_i^{(k)}$ are filter parameters. By stacking convolutional layers, the model captures hierarchical temporal abstractions ranging from local spikes to broader motifs.

CNN-based anomaly detectors are attractive for three main reasons. First, they are highly parallelizable and often train faster than recurrent architectures. Second, they are effective at learning short-term local signatures such as ramps, oscillations, bursts, or repeating micro-patterns. Third, temporal convolutional networks with dilation can expand the receptive field significantly while retaining computational efficiency \cite{bai2018tcn}. In anomaly detection, CNNs are commonly used in forecasting models, convolutional autoencoders, or hybrid spatio-temporal encoders. DeepAnT is a representative CNN-based forecasting detector that uses prediction error as the anomaly signal \cite{munir2019deepant}. MSCRED extends the idea through a multi-scale convolutional recurrent encoder-decoder for multivariate spatio-temporal anomaly detection \cite{zhang2019mscred}.

CNNs are also closely connected to deployment efficiency. Architectures such as MobileNet show that convolutional feature extraction can be made lightweight through depthwise separable convolutions, reducing computation and model size while preserving much of the representational benefit \cite{howard2017mobilenets}. This is highly relevant to edge and industrial anomaly detection, where monitoring systems may operate under strict hardware constraints \cite{antonini2023tinyml}. For a thesis concerned with deployability, CNN literature therefore matters not only as a modeling family, but also as a reference point for efficient neural design.

However, CNN-based detectors also have limitations. Their strongest inductive bias is local rather than long-range. Although dilation and depth increase the effective context window, fixed convolutional kernels may still struggle when anomaly evidence depends on distant events, prolonged dependencies, or regime-level transitions. Moreover, CNN scores usually provide limited semantic explanation. They can indicate that a pattern is statistically unusual, but not why that pattern violates market logic or physical constraints.

\subsection{LSTM-Based Time-Series Anomaly Detectors}

Long Short-Term Memory networks were proposed to address the vanishing-gradient limitations of classical recurrent neural networks and to better capture long-term sequence dependencies \cite{hochreiter1997long}. Their gating mechanism can be written as
\begin{align}
	\mathbf{i}_t &= \sigma(W_i \mathbf{x}_t + U_i \mathbf{h}_{t-1} + \mathbf{b}_i), \\
	\mathbf{f}_t &= \sigma(W_f \mathbf{x}_t + U_f \mathbf{h}_{t-1} + \mathbf{b}_f), \\
	\mathbf{o}_t &= \sigma(W_o \mathbf{x}_t + U_o \mathbf{h}_{t-1} + \mathbf{b}_o), \\
	\tilde{\mathbf{c}}_t &= \tanh(W_c \mathbf{x}_t + U_c \mathbf{h}_{t-1} + \mathbf{b}_c), \\
	\mathbf{c}_t &= \mathbf{f}_t \odot \mathbf{c}_{t-1} + \mathbf{i}_t \odot \tilde{\mathbf{c}}_t, \\
	\mathbf{h}_t &= \mathbf{o}_t \odot \tanh(\mathbf{c}_t).
	\label{eq:sota_lstm}
\end{align}
This structure enables the network to selectively remember, forget, and expose information over time, which made LSTMs a dominant model family in sequence anomaly detection for many years.

LSTM-based anomaly detectors appear in three major forms. First, forecasting-based LSTMs predict future values and use forecast residuals as anomaly scores. Second, encoder-decoder LSTMs reconstruct sequences and identify anomalies through reconstruction mismatch. Third, recurrent probabilistic models combine LSTM dynamics with stochastic latent variables. The LSTM encoder-decoder of Malhotra \emph{et al.} is a foundational example of sequence reconstruction for multi-sensor anomaly detection \cite{malhotra2016encdecad}. Hundman \emph{et al.} showed that LSTM forecasting combined with dynamic thresholding can detect anomalies in spacecraft telemetry under operational constraints \cite{hundman2018spacecraft}. OmniAnomaly integrated gated recurrent modeling with variational inference for multivariate anomaly detection under seasonal behaviour \cite{su2019omnianomaly}.

The key strength of LSTMs is their direct handling of temporal order. Compared with simple CNNs, they are generally better at modeling long sequential dependencies when anomaly evidence depends on accumulated historical context. This explains why they became highly influential in industrial telemetry, cyber-physical monitoring, and energy forecasting tasks. Nevertheless, important limitations remain. Training and inference are less parallelizable than in convolutional architectures because recurrence is sequential. Long-horizon dependency modeling still degrades in practice, despite the gating mechanism. Most importantly for this thesis, LSTM-based detectors usually remain black-box statistical detectors. They can identify that a sequence does not fit learned dynamics, but not whether the detected deviation is physically valid, contextually explainable, or actionable for downstream market decisions.

\subsection{Transformer-Based Time-Series Anomaly Detectors}

Transformers replaced recurrence with self-attention, enabling each time step to interact directly with many others in the sequence \cite{vaswani2017attention}. The scaled dot-product attention mechanism is
\begin{equation}
	\mathrm{Attention}(Q,K,V) = \mathrm{softmax}\!\left(\frac{QK^{\top}}{\sqrt{d_k}}\right)V,
	\label{eq:sota_attention}
\end{equation}
where $Q$, $K$, and $V$ are query, key, and value matrices. In time-series modeling, this allows the detector to capture long-range dependencies, variable periodicities, and complex cross-time associations without explicit recurrence.

Transformer-based anomaly detectors are widely regarded as part of the current state of the art in multivariate time-series anomaly detection. Their main advantage is flexibility: they can learn both local and global dependencies, and they can model cross-channel interactions in a unified attention mechanism. Anomaly Transformer is a representative example, arguing that anomalies manifest through abnormal association discrepancy in the attention structure itself \cite{xu2022anomalytransformer}. TranAD combines Transformer representations with reconstruction-style learning and adversarial components, achieving strong performance across benchmark datasets \cite{tuli2022tranad}. These methods are particularly compelling when anomalies arise not merely from scalar deviations but from disrupted temporal relationships.

Despite these strengths, Transformer-based detectors also introduce trade-offs. Self-attention can be computationally expensive, often scaling quadratically with sequence length. This becomes problematic for long-horizon monitoring or resource-constrained environments. In addition, although attention maps are sometimes interpreted as a form of interpretability, they do not automatically provide domain-level explanations. They may indicate which time steps influenced the model, but they do not directly explain whether the anomaly violates a physical law, a market rule, or an operational constraint. As a result, even sophisticated Transformer-based detectors may remain statistically powerful but operationally opaque.

\subsection{Comparative Synthesis for Energy-System Anomaly Detection}

Table~\ref{tab:sota_detector_families} summarizes the major deep-learning families discussed above from the perspective of time-series anomaly detection in operational energy systems.

\begin{table}[H]
	\centering
	\caption{Comparison of major deep-learning families for time-series anomaly detection in operational energy settings.}
	\label{tab:sota_detector_families}
	\small
	\begin{tabular}{|p{2.4cm}|p{3.1cm}|p{2.8cm}|p{3.2cm}|}
		\hline
		\textbf{Family} & \textbf{Typical anomaly mechanism} & \textbf{Main strength} & \textbf{Main limitation} \\
		\hline
		Autoencoder-based & Reconstruction error, latent density, memory mismatch & Unsupervised learning of normal manifolds; strong multivariate compression & May reconstruct anomalies too well; weak operational semantics \\
		\hline
		CNN-based & Forecasting or reconstruction error from convolutional features & Efficient local motif extraction, multi-scale pattern learning, strong parallelism & Less natural handling of very long dependencies; limited explanation \\
		\hline
		LSTM-based & Forecast residuals, sequence reconstruction, latent probabilistic score & Strong sequential modeling and explicit temporal memory & Slower training/inference; black-box behaviour; limited actionability \\
		\hline
		Transformer-based & Attention discrepancy, reconstruction, forecasting, or hybrid score & Strong long-range dependency modeling and flexible cross-time association learning & Higher computational cost; data-hungry; attention alone is not domain explanation \\
		\hline
	\end{tabular}
\end{table}

A clear progression emerges from the literature. Autoencoders established reconstruction-based unsupervised anomaly detection; CNNs improved efficient local feature extraction; LSTMs improved sequential memory and temporal forecasting; and Transformers improved long-range dependency modeling and representation flexibility. However, all these families share an important limitation in practical energy domains: they are primarily optimized to detect statistical irregularity. They are much less explicit about whether a detected anomaly is physically plausible, contextually justified, economically relevant, or actionable for downstream decisions.

\subsection{Implications for the Thesis}

The reviewed model families provide strong mechanisms for learning statistical
regularities in multivariate time series. However, their outputs are generally not
constrained by explicit physical, calendar, or market knowledge. Statistical
irregularity alone therefore does not establish whether an event is contextually
plausible or operationally relevant.

This limitation motivates a complementary reasoning stage that can examine
statistical anomaly candidates using explicit domain knowledge. The symbolic-AI
literature relevant to this role is reviewed in the following section.

\section{Symbolic AI in Energy Systems}
\label{sec:symbolic_ai_energy_systems}

Symbolic Artificial Intelligence represents domain knowledge explicitly through
facts, rules, relations, constraints, ontologies, or formal inference procedures
\cite{russell2021ai,buchanan1984rulebased,baral2003knowledge,sowa2000knowledge}.
This is particularly relevant in engineering systems, where physical constraints,
operational procedures, market rules, and exception cases should remain visible
and verifiable rather than being represented only implicitly within a learned
model.

Symbolic methods have a long history in power-system diagnosis, alarm processing,
restoration, and operator support. This section reviews that development from
classical expert systems to declarative logic, Answer Set Programming, and
neuro-symbolic integration, with particular attention to rule-aware anomaly
interpretation.



\subsection{What Symbolic AI Means in Engineering Contexts}

At its core, symbolic AI is concerned with \emph{knowledge representation} and \emph{inference}. A symbolic system usually contains a set of explicit facts and a set of rules that determine how new conclusions can be derived. If a system observes that a line is disconnected, a breaker has opened, and a relay has operated, then a symbolic reasoner can infer candidate fault states through predefined logical dependencies. This differs from black-box statistical inference because the reasoning path is explicit and inspectable.

A minimal symbolic rule can be written as
\begin{equation}
	\texttt{conclusion} \leftarrow \texttt{condition}_1,\texttt{condition}_2,\dots,\texttt{condition}_n,
	\label{eq:ch32_symbolic_rule}
\end{equation}
meaning that the conclusion follows whenever the listed conditions hold. More expressive symbolic systems extend this with defaults, negation, priorities, and constraints. The central advantage is not only interpretability, but also the ability to inject \emph{trusted prior knowledge} directly into the decision process. In engineering domains, this is often essential because certain relationships are known a priori and should not be rediscovered from data every time a model is retrained.

From the perspective of energy systems, symbolic AI is attractive for four reasons. First, power-system and market behaviour is strongly constrained by explicit structure, such as network laws, protection logic, balancing conditions, and market rules. Second, engineering decisions often require \emph{justification}: the user wants to know not only that an anomaly was flagged, but why it was confirmed or rejected. Third, the domain involves many \emph{exception cases}, such as holidays, outages, curtailment events, or weather regimes, where ordinary patterns temporarily become acceptable. Fourth, some reasoning tasks are naturally relational rather than numerical; topology, switching sequences, alarm dependencies, and event causality are examples. These properties make symbolic AI particularly relevant for electricity systems even in the era of large-scale machine learning.

\subsection{Classical Expert Systems in Power Engineering}

One of the earliest and most influential forms of symbolic AI in engineering was the \emph{expert system}. Expert systems aimed to emulate human expert reasoning in narrow domains by encoding domain knowledge as rules and applying an inference engine to observed facts \cite{buchanan1984rulebased,giarratano2005expert}. A traditional expert-system architecture contains at least five major components:
\begin{enumerate}
	\item a \textbf{knowledge base}, storing domain rules and structured expertise;
	\item a \textbf{working memory}, storing case-specific facts currently known;
	\item an \textbf{inference engine}, applying rules to derive new conclusions;
	\item an \textbf{explanation module}, justifying conclusions in human-readable form;
	\item and often a \textbf{knowledge acquisition interface}, supporting maintenance and updates.
\end{enumerate}

Table~\ref{tab:expert_system_architecture} summarizes this architecture from an energy-systems viewpoint.

\begin{table}[H]
	\centering
	\caption{Classical expert-system architecture interpreted for energy-system decision support.}
	\label{tab:expert_system_architecture}
	\small
	\begin{tabular}{|p{3.5cm}|p{5.5cm}|p{5.5cm}|}
		\hline
		\textbf{Component} & \textbf{General symbolic role} & \textbf{Energy-system interpretation} \\
		\hline
		Knowledge base & Stores explicit rules and domain heuristics & Protection rules, outage logic, alarm interpretation, restoration priorities, operational heuristics \\
		\hline
		Working memory & Stores case-specific facts & Observed breaker states, relay signals, load conditions, topology changes, market alerts \\
		\hline
		Inference engine & Applies rules to derive conclusions & Suggests likely fault sections, restoration actions, anomaly causes, or rule violations \\
		\hline
		Explanation facility & Justifies why a conclusion was reached & Provides operator-readable reasons for a diagnosis, recommendation, or veto \\
		\hline
		Knowledge acquisition & Supports maintenance and revision of rules & Allows engineers to update expert rules as procedures, equipment, or regulations change \\
		\hline
	\end{tabular}
\end{table}

Two major inference styles are particularly important. In \textbf{forward chaining}, the inference engine starts from known facts and repeatedly fires applicable rules until no new conclusions can be derived. This is suitable when data arrive continuously and the goal is to interpret a stream of evidence, such as alarm processing in control rooms. In \textbf{backward chaining}, the system starts from a hypothesis and checks whether available facts support it. This is suitable when testing a suspected diagnosis or verifying whether a particular abnormal state is consistent with observations.

Historically, expert systems were attractive for power engineering because they matched the way many operational tasks were performed by humans. Control-room operators and system engineers routinely rely on structured procedures, known contingencies, and if--then heuristics. Tasks such as fault diagnosis, alarm filtering, protective-relay interpretation, contingency response, and service restoration all contain substantial symbolic structure. In that sense, the early use of expert systems in power systems was not an artificial imposition of AI onto engineering; it was a formalization of reasoning patterns already used by practitioners.

Yet the limitations of classical expert systems also became clear. Knowledge acquisition was difficult, because domain expertise had to be elicited manually and translated into maintainable rules. Rule sets could become brittle and hard to scale. Continuous sensor streams and uncertain values had to be discretized before the rules could be applied. Most importantly, classical expert systems struggled to learn new regularities automatically from raw data. As energy systems became increasingly data-driven, weather-dependent, and multivariate, purely handcrafted symbolic systems became insufficient. Still, the underlying lesson survived: explicit reasoning remains valuable wherever interpretability, rule compliance, and structured diagnosis matter.

\subsection{Knowledge-Based Reasoning Tasks in Energy Systems}

To understand why symbolic AI remained relevant, it is useful to distinguish the types of reasoning tasks that arise in energy systems. These tasks are not all the same, and different symbolic paradigms serve different purposes.

\textbf{Fault diagnosis and event interpretation.} Grid faults, breaker operations, relay activations, and unusual telemetry patterns often require causal interpretation rather than mere classification. Symbolic rules can encode causal chains such as ``if breaker $B$ opens and relay $R$ trips after line current exceeds threshold, then line section $L$ is a candidate fault region.'' This supports explanation-rich diagnostic reasoning.

\textbf{Alarm processing and reduction.} Modern control rooms can generate large streams of alarms, many of which are redundant, dependent, or secondary consequences of a smaller root-cause event. Symbolic reasoning can suppress cascaded alarm noise, cluster related events, and explain which alarms are primary and which are derivative.

\textbf{Service restoration and switching logic.} Restoration after outages is a procedure-rich domain. Switching actions must satisfy topology constraints, safety rules, sequencing requirements, and service-priority logic. Such reasoning is difficult to express as pure statistical prediction, but natural in symbolic form.

\textbf{Constraint checking and compliance reasoning.} Market and grid operations are governed by explicit admissibility conditions. Some states or recommendations are simply not allowed. Symbolic constraints are particularly valuable in these cases because they can directly rule out forbidden configurations rather than merely assigning them low probability.

\textbf{Context-sensitive anomaly interpretation.} This is especially important for the present thesis. Some anomalies are only meaningful under particular weather, calendar, or market contexts. Symbolic reasoning provides a mechanism for saying that an observed deviation is statistically unusual yet operationally explainable.

These reasoning tasks illustrate why symbolic AI is fundamentally different from generic statistical pattern recognition. In many power applications, the challenge is not only to detect something rare, but to determine whether it is logically consistent, procedurally important, and actionable.

\subsection{From Production Rules to Declarative Knowledge Representation}

Although expert systems were influential, later work in symbolic AI increasingly emphasized \emph{declarative knowledge representation}. In a declarative system, the modeler specifies \emph{what} relationships should hold rather than \emph{how} to execute the reasoning procedure. Logic programming is a natural example of this paradigm. A simple clause may be written as
\begin{equation}
	h \leftarrow b_1,b_2,\dots,b_m,
	\label{eq:ch32_logic_programming}
\end{equation}
where the head $h$ is derived when all body atoms hold. This is appealing because the semantics of the rule are clear and independent of an arbitrary execution order \cite{baral2003knowledge}.

Declarative representations have important advantages in energy domains. They support modular rule design, easier maintenance, and clearer explanation. They also work well with relational knowledge, which is essential for topology modeling, switching-state dependencies, and multi-entity event relationships. For example, one can represent that a given alarm depends on a specific equipment state, or that a restoration action requires a prior switching condition to be satisfied. Such knowledge is far more naturally expressed in a declarative relational form than in a continuous vector.

At the same time, declarative systems expose a key challenge: real energy systems often involve incomplete information, defaults, and exceptions. A rule may generally hold, but be overridden by a holiday, an outage, or an alternative operating regime. This means that \emph{monotonic} logic, where adding new facts never invalidates previous conclusions, is often too rigid. The practical reasoning needs of power systems therefore motivate non-monotonic symbolic paradigms.

\subsection{Answer Set Programming: Stable Models, Defaults, and Constraints}

Answer Set Programming is one of the most important modern frameworks for non-monotonic declarative reasoning \cite{gelfond1988stable,brewka2011asp,gebser2012answer,gebser2019multi}. Its foundation is the stable-model semantics introduced by Gelfond and Lifschitz, which allows logic programs to reason with defaults, exceptions, and incomplete knowledge in a principled way \cite{gelfond1988stable}. An ASP rule typically takes the form
\begin{equation}
	h \leftarrow b_1,\dots,b_m,\ \texttt{not}\ c_1,\dots,\texttt{not}\ c_n,
	\label{eq:ch32_asp_rule}
\end{equation}
where \texttt{not} denotes default negation. This means that $h$ can be concluded when the positive conditions hold and there is no evidence for the negated conditions.

This matters deeply for energy systems because many operational judgments are default-based. For example, a strong anomaly candidate might normally be treated as important unless it is explained by a holiday, a nighttime solar condition, or a known procedural exception. ASP provides a concise way to encode such reasoning. In addition, ASP supports integrity constraints of the form
\begin{equation}
	\leftarrow a_1,\dots,a_k,
	\label{eq:ch32_integrity_constraint}
\end{equation}
which state that the simultaneous truth of $a_1,\dots,a_k$ is inadmissible. Constraints are extremely useful in engineering domains because they directly represent impossible or forbidden conditions.

ASP also supports a particularly powerful form of \emph{explainable refinement}. A system may first propose candidate states or anomalies using a statistical model, and then use ASP to retain only those candidates that survive symbolic scrutiny. This is more expressive than a simple post-hoc rules engine because the ASP layer can reason with defaults, exceptions, and interacting constraints rather than executing rules in a single procedural order.

Operationally, modern ASP solvers such as \texttt{clingo} make this practical by grounding symbolic rules into concrete instances and solving for stable models efficiently \cite{gebser2019multi}. This allows symbolic reasoning to be embedded in data pipelines rather than confined to toy examples.

\subsection{Symbolic AI Under Uncertainty: Fuzzy Rules, Ontologies, and Hybrid Knowledge Structures}

A deeper symbolic-AI review should also acknowledge that not all symbolic reasoning in engineering takes the form of crisp if--then rules. Real energy systems involve uncertainty, imprecision, and semantic heterogeneity. As a result, several adjacent symbolic paradigms have been explored.

\textbf{Fuzzy rule-based systems} relax the hard boundaries of classical symbolic logic by allowing graded membership and approximate reasoning \cite{zadeh1965fuzzy,klir1995fuzzy}. Instead of a binary statement such as ``load is high,'' a fuzzy system may represent \emph{degrees} of highness and combine them through fuzzy inference. This is useful when operator concepts are linguistically meaningful but numerically imprecise.

\textbf{Ontologies and semantic models} focus on defining domain concepts and relations in a shared formal vocabulary \cite{gruber1993ontology,sowa2000knowledge}. In smart-grid contexts, such representations are useful for interoperability, event semantics, and structured data integration across heterogeneous systems. Ontologies are not necessarily inference engines by themselves, but they play an important supporting role in knowledge organization and machine-readable domain semantics.

\textbf{Constraint-based symbolic reasoning} emphasizes admissibility conditions, feasibility checks, and rule compliance. In power systems, many decisions are constrained not only by prediction quality but also by safety and procedure. Constraint-based symbolic methods remain highly relevant whenever a state must be ruled out categorically.

These paradigms matter for the present thesis because they show that symbolic AI is not a single narrow technique. It is a family of methods for making domain knowledge explicit. ASP is selected here not because it is the only symbolic method, but because it provides the right balance of expressiveness, defaults, constraint handling, and practical solver support for anomaly-refinement tasks.

\subsection{Why Symbolic AI Still Matters in Weather-Driven Energy Systems}

The growth of deep learning has not removed the need for symbolic reasoning in
weather-driven energy systems. As renewable generation, market behaviour, and
grid operation become more context-dependent, a statistically unusual observation
may still be explainable through weather, calendar, curtailment, or operating-state
information \cite{amin2005toward,fang2012smartgrid,momoh2012smart}.

Symbolic reasoning remains valuable for three related purposes: it can reject
candidates that conflict with known physical or operational conditions, provide
explicit reasons for confirmation or rejection, and help distinguish alerts that
are operationally relevant from those that are merely statistically unusual. These
roles complement rather than replace statistical learning.

\subsection{From Pure Symbolic Systems to Neuro-Symbolic Energy Intelligence}

Symbolic methods alone are not sufficient for anomaly detection in continuous,
noisy, and high-dimensional energy data. Such systems generally require sensor and
market observations to be converted into symbolic facts before reasoning can take
place, while the relevant temporal relationships may be difficult to specify
manually.

A neuro-symbolic design addresses this limitation by assigning complementary roles
to the two paradigms. A learning component models multivariate temporal behaviour
and produces anomaly candidates, while a symbolic component evaluates those
candidates using explicit contextual and domain constraints. This division is
particularly suitable for electricity-market and grid-monitoring applications,
where both statistical sensitivity and domain-level reasoning are required.

\subsection{Expanded View of Symbolic AI in the Present Thesis}

The present thesis adopts this neuro-symbolic pattern by translating statistical
anomaly candidates and selected contextual variables into symbolic facts that are
evaluated using ASP. The symbolic layer represents persistence, calendar context,
physical plausibility, cross-signal support, and explicit veto information
\cite{repo_thesis_grid_anomaly}.

Since the theoretical foundations of ASP were introduced in
Chapter~\ref{sec:Slimemold}, the repository-specific fact generation, rule
families, and solver integration are described in Chapter~\ref{sec:methods}. The
complete ASP rule base is provided in Appendix~\ref{sec:appendix-netlogo}.

\subsection{Comparative Synthesis}

Table~\ref{tab:symbolic_energy_methods_expanded} summarizes the main symbolic paradigms relevant to energy-system intelligence and situates the present thesis among them.

\begin{table}[H]
	\centering
	\caption{Expanded comparison of symbolic AI paradigms relevant to energy systems.}
	\label{tab:symbolic_energy_methods_expanded}
	\small
	\begin{tabular}{|p{2.8cm}|p{3.5cm}|p{3.5cm}|p{4.7cm}|}
		\hline
		\textbf{Paradigm} & \textbf{Knowledge form} & \textbf{Typical energy role} & \textbf{Main limitation} \\
		\hline
		Expert systems & Production rules and heuristics & Fault diagnosis, alarm processing, restoration guidance, operator support & Brittle maintenance, weak learning from raw data, expensive knowledge engineering \\
		\hline
		Classical logic programming & Facts and logical clauses & Topology reasoning, dependency tracking, procedure checking & Limited treatment of defaults and exceptions in dynamic settings \\
		\hline
		Fuzzy rule systems & Linguistic rules with graded membership & Approximate engineering reasoning under imprecision & Less crisp explanation of hard constraints; rule design remains manual \\
		\hline
		Ontologies / semantic models & Formal domain vocabulary and relations & Interoperability, concept alignment, event semantics, structured metadata reasoning & Do not by themselves solve dynamic anomaly interpretation \\
		\hline
		Answer Set Programming & Non-monotonic rules, defaults, constraints, stable models & Anomaly veto, diagnosis, explainable constraint-based filtering, rule-compliant reasoning & Continuous signals must be abstracted into symbolic facts \\
		\hline
		Neuro-symbolic systems & Learned statistical outputs plus symbolic reasoning & Hybrid anomaly detection, explainable monitoring, rule-aware decision support & Integration quality depends on both learned representations and rule quality \\
		\hline
	\end{tabular}
\end{table}

\subsection{Implications for the Thesis Gap}

The reviewed literature reveals complementary strengths. Data-driven methods can
learn complex patterns from noisy multivariate signals, whereas symbolic methods
provide explicit constraints, explanations, and rule compliance. Neither paradigm
alone fully satisfies both requirements.

Their integration therefore addresses a gap that remains unresolved by purely
statistical or purely symbolic approaches. The following section formulates this
gap more specifically in terms of interpretability and actionability.

\section{The Gap: Interpretability and Actionability in Existing Models}
\label{sec:gap_analysis}

The preceding review shows that deep models provide strong statistical pattern
recognition, while symbolic methods provide explicit knowledge representation and
rule-based reasoning. However, these capabilities are rarely integrated fully in
energy-oriented anomaly detection. The resulting gap is the limited ability of
existing systems to determine whether statistically unusual events are also
interpretable, contextually plausible, and relevant to downstream operational or
market decisions.



\subsection{From Statistical Detection to Usable Intelligence}

As described in Section~\ref{sec:sota_anomaly_detection}, most anomaly detectors
map observations to anomaly scores and then to thresholded decisions. This is
sufficient for identifying statistical irregularity, but it does not by itself
establish whether a detected event is suitable for operational use. A useful
anomaly decision should therefore address at least four questions:
\begin{enumerate}
	\item \textbf{Was the event statistically unusual?}
	\item \textbf{Is the event physically or contextually plausible?}
	\item \textbf{Can the event be explained to an analyst or operator?}
	\item \textbf{Does the event matter for a downstream decision or utility function?}
\end{enumerate}

Most anomaly-detection research focuses primarily on statistical unusualness,
while explanation and downstream relevance are addressed less systematically
\cite{chandola2009anomaly,blazquez2021review}. In energy systems, this limitation
is amplified because the meaning of a deviation depends on weather, calendar
conditions, operating state, and relationships among signals. The unresolved
problem is therefore not only improving anomaly scores, but connecting detection
to understanding and decision relevance.

\subsection{The Interpretability Gap in Black-Box Deep Anomaly Detection}

A major limitation of recent deep anomaly detectors is their restricted interpretability. The literature on interpretable machine learning repeatedly emphasizes that black-box models can achieve high predictive performance while remaining difficult to understand, validate, and trust \cite{doshi2017rigorous,lipton2018mythos,guidotti2018survey,arrieta2020xai,rudin2019stop}. This challenge is especially acute in anomaly detection because anomaly decisions are already ambiguous by nature. If the detector is itself opaque, then it becomes difficult to determine whether a flagged event reflects a meaningful system irregularity, a benign context shift, a data-quality artifact, or a model failure.

This problem affects all major deep-learning families reviewed earlier. Autoencoder-based models often identify anomalies through reconstruction error, but reconstruction mismatch alone does not explain why the pattern is abnormal in domain terms. CNN-based detectors can identify abnormal local motifs, but usually provide little insight into whether those motifs violate physical or market logic. LSTM-based detectors often rely on forecast residuals or latent sequence reconstruction, yet they generally do not reveal which domain rules justify treating a deviation as important. Transformer-based detectors are sometimes presented as more interpretable because they expose attention weights or association patterns, but attention is not the same as explanation; it does not directly state whether an event is plausible, relevant, or consistent with domain knowledge \cite{lipton2018mythos,arrieta2020xai}.

From the perspective of electricity-market and grid monitoring, this interpretability gap is not a minor inconvenience. It affects model validation, operator trust, and deployment viability. In a weather-driven energy system, analysts need to know whether an anomaly was due to absent solar radiation, an unusual wind ramp, a holiday effect, or a suspicious price-load inconsistency. A purely black-box detector rarely provides such explanations explicitly. At best, post-hoc explainability tools may indicate which variables influenced the score. But this is still different from \emph{domain-grounded interpretability}, where the anomaly is expressed in terms of rules, constraints, and operational meaning.

The consequence is that many state-of-the-art anomaly detectors remain \emph{descriptively strong but semantically thin}. They can tell the user that something unusual happened, but not necessarily whether that unusual event should be believed, investigated, ignored, or acted upon. This interpretability gap is one of the main reasons why deep anomaly detection has not fully solved the needs of operational energy analytics.

\subsection{The Actionability Gap in Energy Trading and Grid Operations}

Interpretability alone is not the whole problem. Even if a model produces a plausible and partially explainable anomaly score, a further gap remains between anomaly detection and \emph{actionability}. In practical energy systems, the purpose of detecting anomalies is rarely detection for its own sake. Anomalies are useful because they may indicate operational stress, market opportunity, data inconsistency, control risk, or the need for human investigation. In other words, anomaly detection is often a means to support a downstream decision.

This is especially important in electricity markets. A day-ahead price anomaly, load anomaly, or renewable-generation irregularity has value only insofar as it improves decision support, risk awareness, or strategic positioning. Yet most anomaly-detection papers evaluate performance primarily through statistical criteria rather than decision-relevant outcomes. Similarly, much of the electricity-price forecasting literature emphasizes prediction accuracy while acknowledging that strong statistical forecasting does not automatically translate into strong economic value \cite{weron2014electricity,lago2021forecasting}. The same logic applies to anomaly detection. A model may detect unusual events accurately in a statistical sense but still fail to improve market-oriented decisions if it does not distinguish noise from economically relevant irregularities.

This disconnect can be expressed formally. Suppose a detector produces a binary anomaly signal $z_t$, and suppose an operational or trading action $a_t$ depends on that signal. Then the true downstream goal is not merely to optimize the anomaly score, but to optimize a utility functional
\begin{equation}
	U = \sum_{t=1}^{T} u(a_t, \mathbf{y}_t, z_t),
	\label{eq:gap_utility}
\end{equation}
where $u(\cdot)$ represents task-specific reward, cost, penalty, or benefit. Existing anomaly-detection models typically optimize a surrogate objective related to reconstruction error, forecasting loss, or representation discrepancy, but they rarely close the loop to the utility level. As a result, the literature contains a structural gap between \emph{what is optimized during modeling} and \emph{what matters during decision support}.

The present thesis addresses this issue by evaluating anomaly outputs through a
finance-aware utility formulation rather than stopping at statistical detection
metrics. The evaluation design is presented in Chapter~\ref{sec:implementation},
and the corresponding results are reported in Chapter~\ref{sec:results}.

\subsection{Why Existing Symbolic Approaches Alone Do Not Fully Close the Gap}

Symbolic methods improve explanation, rule enforcement, and procedural reasoning,
but they do not naturally learn subtle nonlinear patterns directly from large,
continuous, and noisy multivariate time series. Deep models provide this statistical
sensitivity, but generally offer less explicit domain-level reasoning.

The gap therefore cannot be closed by selecting one paradigm and excluding the
other. It requires a hybrid architecture in which a learning component identifies
candidate irregularities and a symbolic component evaluates them against explicit
contextual and domain constraints.

\subsection{A Comparative Statement of the Gap}

The literature gap can therefore be summarized more precisely than simply stating that ``existing models lack interpretability.'' A more accurate statement is that existing methods often optimize one or two desirable properties while leaving the others underdeveloped. Table~\ref{tab:gap_comparison} synthesizes this issue.

\begin{table}[H]
	\centering
	\caption{Comparative view of the literature gap in anomaly detection for energy-oriented applications.}
	\label{tab:gap_comparison}
	\small
	\begin{tabular}{|p{2.5cm}|p{2.8cm}|p{2.8cm}|p{2.6cm}|p{3.8cm}|}
		\hline
		\textbf{Paradigm} & \textbf{Statistical sensitivity} & \textbf{Interpretability} & \textbf{Rule consistency} & \textbf{Decision utility / actionability} \\
		\hline
		Black-box deep detectors & High & Low to moderate & Low & Usually implicit or absent \\
		\hline
		Classical symbolic systems & Low on raw continuous data & High & High & Moderate, but depends on handcrafted knowledge coverage \\
		\hline
		Hybrid heuristic pipelines & Moderate & Moderate & Moderate & Often domain-specific and weakly generalized \\
		\hline
		Desired neuro-symbolic framework & High & High & High & Explicitly modeled \\
		\hline
	\end{tabular}
\end{table}

The key observation is that the literature contains many strong components, but relatively few frameworks that combine them in a coherent way for modern energy-anomaly detection. The gap is therefore not simply the absence of one more high-performing detector. It is the absence of a sufficiently integrated framework that is simultaneously:
\begin{itemize}
	\item accurate enough to detect subtle multivariate irregularities,
	\item explicit enough to justify why anomalies are confirmed or rejected,
	\item rule-aware enough to enforce physical and market plausibility,
	\item and decision-aware enough to map anomaly outputs into downstream operational or financial usefulness.
\end{itemize}

\subsection{Research Gap Addressed by the Present Thesis}

The present thesis is positioned precisely at this intersection. Its central claim is that anomaly detection for weather-dependent electricity markets and grid operations should not terminate at statistical irregularity. Instead, it should proceed through a sequence of increasingly meaningful representations:
\begin{equation}
	\begin{aligned}
		\text{raw multivariate signal}
		&\rightarrow \text{forecast residual / anomaly candidate} \\
		&\rightarrow \text{symbolically refined anomaly}
		\rightarrow \text{utility-relevant event}
	\end{aligned}
	\label{eq:gap_pipeline}
\end{equation}
This sequence captures the conceptual contribution of the proposed dCeNN--ELM--ASP framework. The learning layer provides efficient multivariate forecasting and residual-based anomaly detection. The symbolic layer provides explicit plausibility checks, persistence filtering, and co-anomaly reasoning. The utility layer provides a finance-aware interpretation of whether the refined anomaly signals matter for downstream decisions.

This positioning is also directly reflected in the repository implementation. The current codebase includes a strong black-box sequential baseline in the form of an LSTM benchmark \cite{repo_thesis_grid_anomaly}, an explicit ASP-based refinement stage that produces confirmed and vetoed anomalies with reasons \cite{repo_thesis_grid_anomaly}, and a finance-aware mapping layer that evaluates anomaly outputs in cumulative utility terms rather than only statistical terms \cite{repo_thesis_grid_anomaly}. Together, these components define a gap-driven thesis contribution: not merely a new detector, but a more complete anomaly-intelligence framework.

\subsection{Closing Synthesis}

The literature review therefore leads to a clear and concrete gap statement. Deep anomaly detectors have advanced the state of the art in statistical detection, but remain limited in rule-grounded interpretability and decision-level actionability. Symbolic AI methods provide interpretability and rule compliance, but are insufficient alone for learning subtle patterns from continuous multivariate signals. In energy-market and grid settings, where anomalies are contextual, weather-dependent, and operationally consequential, neither paradigm alone is adequate.

The core research gap addressed by this thesis is thus the lack of an integrated, lightweight, neuro-symbolic anomaly-detection framework that can learn from multivariate temporal data, refine anomaly candidates with explicit domain logic, and evaluate their downstream usefulness in action-oriented terms. This gap provides the immediate motivation for the methodology chapter that follows.



